% Options for packages loaded elsewhere
\PassOptionsToPackage{unicode}{hyperref}
\PassOptionsToPackage{hyphens}{url}
%
\documentclass[
]{article}
\usepackage{amsmath,amssymb}
\usepackage{iftex}
\ifPDFTeX
  \usepackage[T1]{fontenc}
  \usepackage[utf8]{inputenc}
  \usepackage{textcomp} % provide euro and other symbols
\else % if luatex or xetex
  \usepackage{unicode-math} % this also loads fontspec
  \defaultfontfeatures{Scale=MatchLowercase}
  \defaultfontfeatures[\rmfamily]{Ligatures=TeX,Scale=1}
\fi
\usepackage{lmodern}
\ifPDFTeX\else
  % xetex/luatex font selection
\fi
% Use upquote if available, for straight quotes in verbatim environments
\IfFileExists{upquote.sty}{\usepackage{upquote}}{}
\IfFileExists{microtype.sty}{% use microtype if available
  \usepackage[]{microtype}
  \UseMicrotypeSet[protrusion]{basicmath} % disable protrusion for tt fonts
}{}
\makeatletter
\@ifundefined{KOMAClassName}{% if non-KOMA class
  \IfFileExists{parskip.sty}{%
    \usepackage{parskip}
  }{% else
    \setlength{\parindent}{0pt}
    \setlength{\parskip}{6pt plus 2pt minus 1pt}}
}{% if KOMA class
  \KOMAoptions{parskip=half}}
\makeatother
\usepackage{xcolor}
\usepackage[margin=2cm]{geometry}
\usepackage{longtable,booktabs,array}
\usepackage{calc} % for calculating minipage widths
% Correct order of tables after \paragraph or \subparagraph
\usepackage{etoolbox}
\makeatletter
\patchcmd\longtable{\par}{\if@noskipsec\mbox{}\fi\par}{}{}
\makeatother
% Allow footnotes in longtable head/foot
\IfFileExists{footnotehyper.sty}{\usepackage{footnotehyper}}{\usepackage{footnote}}
\makesavenoteenv{longtable}
\usepackage{graphicx}
\makeatletter
\def\maxwidth{\ifdim\Gin@nat@width>\linewidth\linewidth\else\Gin@nat@width\fi}
\def\maxheight{\ifdim\Gin@nat@height>\textheight\textheight\else\Gin@nat@height\fi}
\makeatother
% Scale images if necessary, so that they will not overflow the page
% margins by default, and it is still possible to overwrite the defaults
% using explicit options in \includegraphics[width, height, ...]{}
\setkeys{Gin}{width=\maxwidth,height=\maxheight,keepaspectratio}
% Set default figure placement to htbp
\makeatletter
\def\fps@figure{htbp}
\makeatother
\setlength{\emergencystretch}{3em} % prevent overfull lines
\providecommand{\tightlist}{%
  \setlength{\itemsep}{0pt}\setlength{\parskip}{0pt}}
\setcounter{secnumdepth}{-\maxdimen} % remove section numbering
\ifLuaTeX
\usepackage[bidi=basic]{babel}
\else
\usepackage[bidi=default]{babel}
\fi
\babelprovide[main,import]{vietnamese}
% get rid of language-specific shorthands (see #6817):
\let\LanguageShortHands\languageshorthands
\def\languageshorthands#1{}
\ifLuaTeX
  \usepackage{selnolig}  % disable illegal ligatures
\fi
\usepackage{bookmark}
\IfFileExists{xurl.sty}{\usepackage{xurl}}{} % add URL line breaks if available
\urlstyle{same}
\hypersetup{
  pdflang={vi-VN},
  hidelinks,
  pdfcreator={LaTeX via pandoc}}

\author{}
\date{}

\begin{document}

\subsection{1. Yêu cầu chức năng (Functional Requirements - FR)}\label{yuxeau-cux1ea7u-chux1ee9c-nux103ng-functional-requirements---fr}

\begin{longtable}[]{@{}
  >{\raggedright\arraybackslash}p{(\columnwidth - 6\tabcolsep) * \real{0.2500}}
  >{\raggedright\arraybackslash}p{(\columnwidth - 6\tabcolsep) * \real{0.2500}}
  >{\raggedright\arraybackslash}p{(\columnwidth - 6\tabcolsep) * \real{0.2500}}
  >{\raggedright\arraybackslash}p{(\columnwidth - 6\tabcolsep) * \real{0.2500}}@{}}
\toprule\noalign{}
\endhead
\bottomrule\noalign{}
\endlastfoot
\textbf{Mã ID} & \textbf{Tên Yêu Cầu} & \textbf{Mô tả chi tiết} & \textbf{MoSCoW} \\
FR-01 & Xem danh mục đồ uống & Hiển thị danh sách sản phẩm, danh mục, giá, tồn kho (stock) và trạng thái (active/inactive). & Must \\
FR-02 & Tùy chỉnh món & Cho phép chọn size (S/M/L), số lượng (qty \textgreater= 1) và thêm các loại topping đi kèm. & Must \\
FR-03 & Đăng ký tài khoản & Tạo tài khoản với các ràng buộc chuẩn hóa: email (chỉ nhận các domain phổ biến/edu), số điện thoại chuẩn VN và họ tên định dạng chữ. & Must \\
FR-04 & Đăng nhập \& Xác thực & Đăng nhập bằng email, mật khẩu. Hệ thống phân quyền dựa trên user\_role (customer, barista). & Must \\
FR-05 & Tạo đơn hàng (Checkout) & Khởi tạo đơn hàng từ giỏ với trạng thái PENDING. Hệ thống tự động kiểm tra số lượng tồn kho (stock) của sản phẩm trước khi tạo. & Must \\
FR-06 & Dọn dẹp đơn hàng tự động & Tự động hủy (CANCELLED) các đơn hàng PENDING bị quá hạn hoặc các đơn PENDING cũ của người dùng, đồng thời hoàn trả lại tồn kho. & Must \\
FR-07 & Thanh toán trực tuyến & Thanh toán qua cổng PayOS hoặc Ví MoMo, xác thực thông qua Webhook/IPN và chuyển trạng thái đơn sang PAID. & Must \\
FR-08 & Lịch sử \& Chi tiết đơn & Khách hàng chỉ được xem đơn của mình. Nhân viên/Barista có quyền xem toàn bộ đơn hàng trong hệ thống. & Must \\
FR-09 & Quản lý trạng thái đơn (Barista) & Barista tiếp nhận đơn PAID và lần lượt cập nhật trạng thái hợp lệ: PREPARING -\textgreater{} READY -\textgreater{} COMPLETED. & Must \\
FR-10 & Tích lũy điểm thưởng & Tự động cộng điểm loyaltyPoints (1 điểm / 10.000đ) vào tài khoản khách hàng khi thanh toán đơn thành công. & Should \\
FR-11 & Giảm giá theo hạng thành viên & Tự động tính toán và áp dụng giảm giá (5\%, 10\%, 15\%) tùy theo hạng (Bạc, Vàng, Kim Cương) dựa trên điểm tích lũy khi tạo đơn. & Should \\
\end{longtable}

\subsection{}\label{section}

\subsection{2. Yêu cầu phi chức năng (Non-Functional Requirements - NFR)}\label{yuxeau-cux1ea7u-phi-chux1ee9c-nux103ng-non-functional-requirements---nfr}

\begin{longtable}[]{@{}
  >{\raggedright\arraybackslash}p{(\columnwidth - 6\tabcolsep) * \real{0.2484}}
  >{\raggedright\arraybackslash}p{(\columnwidth - 6\tabcolsep) * \real{0.2484}}
  >{\raggedright\arraybackslash}p{(\columnwidth - 6\tabcolsep) * \real{0.2548}}
  >{\raggedright\arraybackslash}p{(\columnwidth - 6\tabcolsep) * \real{0.2484}}@{}}
\toprule\noalign{}
\endhead
\bottomrule\noalign{}
\endlastfoot
\textbf{Mã ID} & \textbf{Thuộc tính} & \textbf{Mô tả chi tiết (Dựa trên code Backend)} & \textbf{MoSCoW} \\
NFR-01 & Toàn vẹn dữ liệu & Hệ thống phải sử dụng Transaction và Optimistic Locking (thông qua trường version của Product) để ngăn chặn lỗi quá tải tồn kho khi nhiều người cùng đặt một món. & Must \\
NFR-02 & Hiệu năng & Thời gian phản hồi API (đặc biệt là GET /products) phải tối ưu (\textless{} 500ms). Hệ thống xử lý mượt mà khi nhiều khách đặt cùng lúc thông qua cơ chế Optimistic Locking để không bị lỗi trừ âm tồn kho. & Must \\
NFR-03 & Tính khả dụng & Frontend sử dụng Next.js thiết kế theo hướng mobile-first/PWA, có khả năng duy trì giỏ hàng khi tải lại trang. trạng thái đơn hàng phải tuân thủ nghiêm ngặt luồng chuyển đổi một chiều. Có cơ chế retry (thử lại) tự động 3 lần ở backend nếu xảy ra xung đột khi tạo đơn. & Must \\
NFR-04 & Bảo mật & Mật khẩu băm bằng bcrypt, bảo vệ API bằng JWT, phải validate chặt chẽ dữ liệu đầu vào (ví dụ:

chống SQL Injection/XSS ngầm qua Regex). & Must \\
NFR-05 & Khả năng triển khai & Hệ thống (UI, API, DB) phải được đóng gói bằng Docker. Tích hợp pipeline CI/CD chạy tự động qua các bước (Lint -\textgreater{} Test -\textgreater{} Build -\textgreater{} Deploy) và triển khai lên môi trường Production (Render/Railway...). & Must \\
\end{longtable}

\textbf{3.Tiêu Chí Tiếp Nhận:}

\textbf{3.1.Yêu cầu ch}ức năng \textbf{(Functional Requirements - FR):}

\begin{longtable}[]{@{}
  >{\raggedright\arraybackslash}p{(\columnwidth - 2\tabcolsep) * \real{0.5000}}
  >{\raggedright\arraybackslash}p{(\columnwidth - 2\tabcolsep) * \real{0.5000}}@{}}
\toprule\noalign{}
\endhead
\bottomrule\noalign{}
\endlastfoot
\textbf{Mã ID} & Tiêu Chí Tiếp Nhận \\
FR-01 & Khi khách hàng truy cập menu, hệ thống hiển thị danh sách đồ uống cùng tên, giá, danh mục, tồn kho và trạng thái. Sản phẩm inactive không được phép đặt hàng. \\
FR-02 & Khách hàng có thể chọn size S/M/L, số lượng từ 1 trở lên và topping. Tổng tiền được tính đúng theo lựa chọn. \\
FR-03 & Đăng ký thành công khi email, số điện thoại và họ tên hợp lệ, email chưa tồn tại. Hệ thống từ chối dữ liệu không hợp lệ. \\
FR-04 & Người dùng đăng nhập thành công với email và mật khẩu chính xác. Hệ thống từ chối thông tin sai và phân quyền đúng theo vai trò customer hoặc barista. \\
FR-05 & Khi đặt hàng hợp lệ và tồn kho đủ, hệ thống tạo đơn hàng với trạng thái PENDING. Nếu tồn kho không đủ, hệ thống từ chối tạo đơn. \\
FR-06 & Đơn PENDING quá thời hạn hoặc thuộc trường hợp cần dọn dẹp phải chuyển sang CANCELLED và hoàn trả đúng số lượng tồn kho. \\
FR-07 & Khi cổng thanh toán xác nhận giao dịch thành công qua Webhook/IPN hợp lệ, đơn hàng chuyển sang PAID. Giao dịch thất bại không được ghi nhận là đã thanh toán. \\
FR-08 & Khách hàng chỉ xem được đơn hàng thuộc tài khoản của mình. Barista có quyền xem danh sách đơn hàng theo quyền được cấp. \\
FR-09 & Barista có thể cập nhật đơn theo trình tự PREPARING → READY → COMPLETED. Hệ thống từ chối chuyển trạng thái không hợp lệ. \\
FR-10 & Khi đơn thanh toán thành công, hệ thống cộng điểm theo tỷ lệ 1 điểm/10.000 đồng. Không cộng điểm trùng cho cùng một đơn. \\
FR-11 & Hệ thống áp dụng đúng mức giảm giá theo hạng thành viên và điểm tích lũy tương ứng \\
\end{longtable}

\textbf{3.2.Yêu cầu phi chức năng (Non-Functional Requirements - NFR):}

\begin{longtable}[]{@{}
  >{\raggedright\arraybackslash}p{(\columnwidth - 2\tabcolsep) * \real{0.5000}}
  >{\raggedright\arraybackslash}p{(\columnwidth - 2\tabcolsep) * \real{0.5000}}@{}}
\toprule\noalign{}
\endhead
\bottomrule\noalign{}
\endlastfoot
\textbf{Mã ID} & \textbf{Tiêu Chí Chấp Nhận} \\
NFR-01 & Khi nhiều khách hàng đặt cùng một sản phẩm, tồn kho không bị trừ vượt số lượng thực tế; dữ liệu đơn hàng và tồn kho phải nhất quán khi giao dịch thất bại. \\
NFR-02 & API GET /products có thời gian phản hồi dưới 500ms trong môi trường và điều kiện kiểm thử được quy định. \\
NFR-03 & Giao diện hiển thị và sử dụng được trên điện thoại; giỏ hàng vẫn được duy trì sau khi tải lại trang; hệ thống tự thử lại tối đa 3 lần khi xảy ra xung đột tạo đơn. \\
NFR-04 & Mật khẩu trong cơ sở dữ liệu không được lưu dưới dạng văn bản thuần; API được bảo vệ bằng JWT; dữ liệu đầu vào không hợp lệ phải bị từ chối. \\
NFR-05 & Frontend, Backend và Database có thể khởi chạy theo cấu hình Docker; pipeline CI/CD thực hiện thành công các bước Lint, Test, Build và Deploy. \\
\end{longtable}

\begin{enumerate}
\def\labelenumi{\arabic{enumi}.}
\setcounter{enumi}{3}
\item
  \textbf{sơ đồ use-case tổng quát:}

  \includegraphics[width=6.82431in,height=6.25in]{media/image1.png}
\end{enumerate}

\end{document}
